\ssr{Введение}

Цели данной лабораторной работы:
\begin{enumerate}
	\item изучение механизмов синхронизации потоков с использованием сигналов, семафоров, мьютексов, условных переменных и барьеров;
	\item приобретение практических навыков применения средств синхронизации в многопоточных приложениях.
\end{enumerate}

\ssr{Выполнение}

\ssection{Задание 1}

Семафоры (struct sembuf) -- переменная определенного типа, которая доступна параллельным процессам для проведения над ней только двух операций: инкремент P и декремент V.

Мьютексы (pthread\_mutex\_t) -- по сути бинарный семафор, т.е. позволяет разделять доступ к ресурсу нескольким потокам.

Сигналы (int sig, struct sigaction) -- небольшие положительные целые числа, получение процессом которого выглядит как возникновение прерывания: процесс прерывает исполнение, и управление передается функции-обработчику сигнала. 

Условные переменные (pthread\_cond\_t) -- переменная типа, которая обеспечивает блокирование одного или нескольких потоков до тех пор, пока не поступит сигнал, или не пройдет максимально установленное время ожидания.

Барьеры (pthread\_barrier\_t) -- ограничитель потоков, при инициализации которого задается, какое количество потоков он должен остановить; потоки, преодолевающие барьер (дошедшие до него), не могут продолжать свое выполнение пока на этот барьер не наберет указанное при инициализации количество потоков.

\ssection{Задание 2}

Выбранные методы синхронизации: семафоры, мьютексы, сигналы.

%1. Задачи для каждого метода синхронизации должны быть различными.
%2. Задачи должны наглядно демонстрировать выбранный метод
%синхронизации и учитывать особенности его применения.

%\lstinputlisting[style=cstyle, caption={Семафоры}]{sem.c}

%\lstinputlisting[style=cstyle, caption={Мьютексы}]{mutex.c}

%\lstinputlisting[style=cstyle, caption={Сигналы}]{sig.c}

%\begin{figure}[!htbp]
%	\centering
%	%\includegraphics[width=0.6\linewidth]{""}
%	\caption{}
%\end{figure}
%\FloatBarrier

\ssection{Дополнительное задание}

%Напишите программу, решающую задачу синхронизации потоков различных
%не родственных процессов. Средство/а синхронизации выберите, исходя из
%особенностей задачи и реализации

%\lstinputlisting[style=cstyle, caption={}]{}

\ssr{Вывод}

В ходе лабораторной работы были:  
\begin{itemize}
	\item изучены механизмы синхронизации потоков с использованием сигналов, семафоров, мьютексов, условных переменных и барьеров;
	\item приобретены практические навыки применения средств синхронизации в многопоточных приложениях.
\end{itemize}

\ssr{Контрольные вопросы}

\begin{enumerate}
	\item \textit{В чем разница между мьютексом и бинарным семафором?}\newline 
	
	\item \textit{Объясните, что такое «состояние гонки» (race condition) и как мьютекс помогает его избежать.}\newline 
	
	\item \textit{Что такое «взаимная блокировка» (deadlock)? Приведите пример фрагмента кода с использование мьютексов, который может привести к «взаимной блокировке».}\newline 
	
	\begin{lstlisting}[style=cstyle]
	//
	\end{lstlisting}
	\item \textit{Почему условная переменная используется в паре с мьютексом? Что произойдет, если убрать мьютекс?}\newline 
	
	\item \textit{Каков жизненный цикл семафора System V (semget, semop, semctl) и в чем его отличие от семафора POSIX (именованного/неименованного)?}\newline 
	
	\item \textit{В чем опасность использования глобальных переменных в обработчиках сигналов?}\newline 
	
	\item \textit{Возможно ли с помощью функции kill() отправить сигнал конкретному	потоку процесса?}\newline 
	Да, так как ее сигнатура имеет вид int kill(pid\_t pid, int sig).
	\item \textit{Каким значением можно проинициализировать семафор?}\newline 
	
	\item \textit{В каких случаях удачным решением для синхронизации потоков будет использование барьера?}\newline 
	
	\item \textit{Как можно передать экземпляр семафора для использования его, например, в задаче читатели-писатели, если потоки-писатели создает родительский процесс, а потоки-читатели - его дочерний процесс. Считать общим буфером текстовый файл на диске.}\newline 
	
\end{enumerate}
